\documentclass[11pt]{article}
\usepackage[margin=1in]{geometry}
\usepackage{amsmath,amssymb,amsthm,mathtools}
\usepackage[T1]{fontenc}
\usepackage{lmodern,microtype}
\usepackage[colorlinks=true,linkcolor=blue,citecolor=blue,urlcolor=blue]{hyperref}
\hypersetup{pdftitle={Infinitely many cyclic Haagerup--Izumi fusion categories},
 pdfsubject={Proof draft and mathematical audit, September 7, 2026}}
\numberwithin{equation}{section}
\newtheorem{theorem}{Theorem}[section]
\newtheorem{proposition}[theorem]{Proposition}
\newtheorem{lemma}[theorem]{Lemma}
\newtheorem{corollary}[theorem]{Corollary}
\theoremstyle{definition}
\newtheorem{definition}[theorem]{Definition}
\theoremstyle{remark}
\newtheorem{remark}[theorem]{Remark}
\DeclareMathOperator{\Res}{Res}
\DeclareMathOperator{\Tr}{Tr}
\DeclareMathOperator{\FPdim}{FPdim}
\DeclareMathOperator{\Gal}{Gal}
\newcommand{\Z}{\mathbb Z}
\newcommand{\Q}{\mathbb Q}
\newcommand{\C}{\mathbb C}
\newcommand{\R}{\mathbb R}
\newcommand{\F}{\mathbb F}
\newcommand{\one}{\mathbf1}
\allowdisplaybreaks[2]
\title{Infinitely many cyclic Haagerup--Izumi fusion categories}
\author{}
\date{September 7, 2026}

\begin{document}
\maketitle
\begin{abstract}
We construct complex spherical fusion categories with cyclic
Haagerup--Izumi fusion rules at every odd order $n\geq3$ for which
$-1$ belongs to the subgroup generated by $2$ in $(\mathbb Z/n\mathbb Z)^\times$.
In particular, the orders $n=2^m+1$, $m\geq1$, give infinitely many
distinct ranks. The construction starts with explicit double-sine samples.
A finite Fourier evaluation proves the quadratic identities; a contour
difference and induction prove the ordered cubic identities. An unramified
convolution lift supplies the exact field automorphism needed to treat
the remaining indices. Algebraic reconstruction and Galois conjugation
then give spherical and pseudo-unitary categorifications, respectively.
\end{abstract}

\noindent\textbf{AI-use disclosure.}
AI systems were used to conduct the numerical searches, including writing
and running computational code, to develop the proof
arguments, and to write this preliminary \LaTeX{} draft.

\section{Statement and structure of the proof}

For $G=\mathbb Z/n\mathbb Z$, the cyclic Haagerup--Izumi based ring
has basis $\alpha^a,\rho_a$ ($a\in G$) and multiplication
\begin{align}
 \alpha^a\alpha^b&=\alpha^{a+b}, &
 \alpha^a\rho_b&=\rho_{a+b}, &
 \rho_a\alpha^b&=\rho_{a-b},\label{intro:fusion1}\\
 \rho_a\rho_b&=\alpha^{a-b}+\sum_{c\in G}\rho_c.
 \label{intro:fusion2}
\end{align}
Its noninvertible Frobenius--Perron dimension is
$d=(n+\sqrt{n^2+4})/2$.

\begin{theorem}\label{intro:main}
Let $n\geq3$ be odd. Suppose that $2^m\equiv-1\pmod n$ for some
$m\geq1$. There exists a complex spherical fusion category with fusion
rules \eqref{intro:fusion1}--\eqref{intro:fusion2} and noninvertible
categorical dimension $\delta=-1/d$. There also exists a pseudo-unitary
categorification of the same based ring. Consequently there are infinitely
many pairwise inequivalent cyclic Haagerup--Izumi fusion categories.
\end{theorem}

The argument below is independent of any conjectural identification of
the coefficient fields with ray class fields or of the samples with Stark
units. All finite sums and polynomial identities in the proof are exact.
Numerical verification is recorded separately and is not used as a proof
of an identity for unbounded $n$.

Throughout, $A,B$ denote positive real periods, while $\mathsf A,\mathsf B$
denote matrices. All matrix and phase indices are in $G$ unless a signed
integer representative is expressly chosen. The Kronecker symbol is
written $\mathbf1_{\{\cdot\}}$, to distinguish it from the number $\delta$.

% Standalone body fragment. Requires amsmath, amssymb, amsthm.
% Environment names used: proposition, lemma, proof.
% Bibliography key used: BSS.

\section{The analytic candidate and its quadratic identities}

Fix an odd integer $n\geq3$, and put
\[
 d=\frac{n+\sqrt{n^2+4}}2,\qquad \delta=-d^{-1},\qquad
 \tau=\delta-1,\qquad A=d+1,\qquad B=1+d^{-1}.
\]
Then
\[
 \delta^2=n\delta+1,\qquad A-B=n,\qquad
 \Omega:=A+B=AB,\qquad A/B=d,\qquad \tau=-B.
\]
Let $\gamma(z)=\gamma^{(2)}(z;A,B)$ be the hyperbolic gamma
function in the normalization of \cite{BSS}. Thus
\begin{align}
 \gamma(z+A)&=2\sin(\pi z/B)\gamma(z),&
 \gamma(z+B)&=2\sin(\pi z/A)\gamma(z),\label{foundation:shifts}\\
 \gamma(z)\gamma(\Omega-z)&=1,&
 \mathop{\rm Res}_{z=0}\gamma(z)&=\frac{\sqrt\Omega}{2\pi}.
 \label{foundation:reflection}
\end{align}
Its poles are $-rA-sB$, $r,s\geq0$; its zeros are
$\Omega+rA+sB$. It is positive on $(0,\Omega)$.
Write $F(z)=\gamma(B+z)^{-1}$. The preceding identities give
\[
 F(0)=\sqrt d,\qquad F(n-z)=F(z)^{-1},\qquad
 F(-a)=\frac{(-1)^{a+1}}{F(a)}\quad(a\in\mathbb Z\setminus\{0\}).
\]
Define real phases, indexed by $\mathbb Z_n$, by
\begin{equation}\label{foundation:phases}
 W_0=-1,\qquad
 W_a=-(-1)^{a(a+n)/2}\frac{F(a)}{\sqrt d}\quad(1\leq a<n).
\end{equation}
In particular $W_aW_{-a}=\delta$ for $a\ne0$.
Define the core and its rescaling by
\begin{equation}\label{foundation:core}
 \mathsf B_{a,b}=
 \begin{cases}
  \delta-2,&a=b=0,\\
  -1,&a=0\text{ or }b=0\text{ or }a=b,\quad(a,b)\ne(0,0),\\
  W_aW_{b-a}/W_b,&\text{otherwise},
 \end{cases}
 \qquad \mathsf A_{a,b}=\tau^{-1}\mathsf B_{a,b}.
\end{equation}
For $b\ne0$ the phase quotient in \eqref{foundation:core} applies
also at $a=0,b$. The inverse norms imply the order-three symmetry
$\mathsf B_{a,b}=\mathsf B_{b-a,-a}$.

\begin{proposition}\label{foundation:quadratics}
The matrix \eqref{foundation:core} satisfies
\begin{equation}\label{foundation:Q}
 \sum_{a\in\mathbb Z_n}\mathsf A_{g+a,h}\mathsf A_{h,a}
 =\mathbf1_{g=0}-\delta^{-1}\mathbf1_{h=0}
 \qquad(g,h\in\mathbb Z_n).
\end{equation}
More precisely, for $0<h<n$, put
\[
 \widehat{\mathsf B}_h(k)=\sum_{a=0}^{n-1}\mathsf B_{a,h}
              e^{-2\pi ika/n},\qquad
 \omega\equiv\frac{2\pi k}{n}+\pi(h-1)\pmod{2\pi},
 \quad -\pi\leq\omega\leq\pi.
\]
Then
\begin{equation}\label{fd:fourier}
 \widehat{\mathsf B}_h(k)
 =-B e^{-ih\omega/2}
   \gamma\left(\frac{\Omega+h}{2}-\frac{\Omega\omega}{2\pi}\right)
   \gamma\left(\frac{\Omega+h}{2}+\frac{\Omega\omega}{2\pi}\right).
\end{equation}
\end{proposition}

\begin{proof}
Set $R_h(z)=F(z)/F(z-h)$ and
$K_h(z)=\pi e^{-i\omega z}R_h(z)/\sin\pi z$.
At the integer samples, including $a=0,h$,
\begin{equation}\label{foundation:sample}
 \mathsf B_{a,h}+B\mathbf1_{a=h}
 =\frac{(-1)^{h(a+1)}}{\sqrt d\,F(h)}R_h(a).
\end{equation}
Put $\alpha=1/A$ and $\beta=1/B$, so $\alpha+\beta=1$.
The period shifts and the elementary identity
\[
 (-1)^h\sin(\pi\alpha z)\sin(\pi\beta(z-h))
 -\sin(\pi\beta z)\sin(\pi\alpha(z-h))
 =\sin(\pi\alpha h)\sin\pi z
\]
give
\begin{equation}\label{foundation:kernel-difference}
 K_h(z+n)-K_h(z)
 =-4\pi\sin(\pi h/A)e^{-i\omega z}\gamma(-z)\gamma(z-h).
\end{equation}

Let $L=\lfloor h/B\rfloor$ and choose
$\sigma=-\varepsilon$, where
$0<\varepsilon<\min\{1,(L+1)B-h\}$. In the strip
$\sigma<\mathop{\rm Re}z<\sigma+n$, the poles of $K_h$ are
$0,\ldots,n-1$ and $p_\ell=h-\ell B$, $1\leq\ell\leq L$.
There are no numerator poles because the first pole of $F$ is at
$A>n$. The irrationality of $B$ excludes coincidences, and
$K_h(z+n)$ is regular at the additional poles.
The standard gamma asymptotic gives, uniformly on this finite strip,
\[
 |R_h(x+it)|=O(e^{-\pi h|t|/\Omega}).
\]
Thus the horizontal sides of the residue rectangle vanish, including
when $|\omega|=\pi$.

Write $p_0=h$ and
\[
 r_\ell=\mathop{\rm Res}_{z=p_\ell}
         e^{-i\omega z}\gamma(-z)\gamma(z-h),\qquad
 I_\sigma=\int_{\sigma-i\infty}^{\sigma+i\infty}
         e^{-i\omega z}\gamma(-z)\gamma(z-h)\frac{dz}{i}.
\]
The rectangle and \eqref{foundation:kernel-difference} give
\begin{equation}\label{foundation:rectangle}
 \sum_{a=0}^{n-1}\mathop{\rm Res}_{z=a}K_h(z)
 =-2\sin(\pi h/A)I_\sigma
  -4\pi\sin(\pi h/A)\sum_{\ell=1}^{L}r_\ell.
\end{equation}
The two-gamma Fourier integral \cite[Eq.~(2.17)]{BSS}, with
parameters $f=-h$, the other shift zero, and
$\alpha_{\rm BSS}=(\Omega+h-\Omega\omega/\pi)/2$, states that
\begin{align}
 I_{\rm sep}
 &=\sqrt\Omega\,e^{-ih\omega/2}\gamma(-h)
   \gamma\left(\frac{\Omega+h}{2}-\frac{\Omega\omega}{2\pi}\right)
   \gamma\left(\frac{\Omega+h}{2}+\frac{\Omega\omega}{2\pi}\right),
   \label{foundation:beta}\\
 I_{\rm sep}&=I_\sigma+2\pi\sum_{\ell=0}^{L}r_\ell.
   \label{foundation:continued-contour}
\end{align}
Here the separating contour is continued meromorphically from negative
$h$ and $|\omega|<\pi$. The second equality lists every pole that
has crossed the vertical line. At the final integer $h$, no opposite
poles collide: $h<n<A$ cannot equal $rA+sB$ with $r,s\geq0$.
The endpoint frequencies follow by continuity in the convergent
positive-$h$ regime.

Substituting \eqref{foundation:continued-contour} into
\eqref{foundation:rectangle} cancels the residues with $\ell\geq1$.
The remaining residue is
\[
 r_0=\frac{\sqrt\Omega}{2\pi}e^{-i\omega h}\gamma(-h).
\]
Now use
\[
 2\sin(\pi h/A)\gamma(-h)=(-1)^hF(h),\qquad
 \frac{\sqrt\Omega}{\sqrt d}=B.
\]
Multiplication of the integer residue sum by
$(-1)^h/(\sqrt d F(h))$ yields the right side of
\eqref{fd:fourier} plus $Be^{-i\omega h}$.
Since $e^{-i\omega h}=e^{-2\pi ikh/n}$, this extra term cancels
exactly the correction in \eqref{foundation:sample}.

For completeness, reflection of the evaluated gamma factors gives
\begin{equation}\label{foundation:spectral-inverse}
 \widehat{\mathsf B}_h(k)\widehat{\mathsf B}_{n-h}(-k)
 =B^2e^{-2\pi ikh/n}.
\end{equation}
Here is the phase check. By exchanging $h,n-h$ and conjugating,
it suffices to take $h$ odd and $0\leq k\leq(n-1)/2$.
Put $\lambda_k=\Omega k/n$. The first gamma arguments are
$x_\pm=(\Omega+h)/2\pm\lambda_k$, and the other two are
$y_1=\Omega+(n-h)/2-\lambda_k$ and
$y_2=(n-h)/2+\lambda_k$. Since
$x_++y_1=\Omega+A$ and $x_-+y_2=A$, their product reduces to
\[
 \frac{\sin\bigl(\pi[(h-n)/(2B)+Ak/n]\bigr)}
      {\sin\bigl(\pi[(\Omega+h)/(2A)-Bk/n]\bigr)}
 =(-1)^{(h-n)/2+k}.
\]
The two arguments divided by $\pi$ sum to $(h-n+2)/2+k$;
the denominator argument lies strictly between zero and one.
The exponential factors give \eqref{foundation:spectral-inverse}.
Finally $\mathsf B_{h,a}=\mathsf B_{a-h,n-h}$, so Fourier inversion
of \eqref{foundation:spectral-inverse} proves
\eqref{foundation:Q} for $h\ne0$. For $h=0$, it follows directly
from the fixed axes and $\delta^2=n\delta+1$.
\end{proof}

\begin{lemma}\label{foundation:algebraicity}
Every phase $W_a$, and every sampled value $F(a)$ with $0\leq a<n$,
is algebraic over $\mathbb Q$.
\end{lemma}

\begin{proof}
This uses only reality, the inverse norms, and one quadratic column.
Put $q=(n+1)/2$, $u_a=\mathsf B_{a,1}$, and
$v_a=\mathsf B_{1,a}$. Then
\[
 u_0=u_1=v_0=v_1=-1,\qquad u_a=u_{1-a},\qquad
 v_a=\delta/u_a\quad(a\ne0,1).
\]
Let $x_j=u_{q+j}$ and $y_j=v_{q+j}$. Both sequences are even.
Their symmetric circulants $T$ and $V$ satisfy
$TV=\tau^2 I$, by \eqref{foundation:Q} at $h=1$.
If $c_j$ are the coefficients of $T^{-1}$, then
$c_{a-q}=v_a/\tau^2$.

On the real affine space $u_a=u_{1-a}$, $u_0=u_1=-1$,
consider, locally where every free coordinate and $\det T$ are nonzero,
\[
 \Phi(u)=\log|\det T|+\eta\sum_{a=2}^{n-1}\log|u_a|,
 \qquad \eta=-\frac{n\delta}{\tau^2}>0.
\]
For a free reflection orbit represented by $a$, its derivative is
the orbit multiplicity times
$nc_{a-q}+\eta/u_a=0$. Thus the given $u$ is critical.
Every directional derivative $H$ of $T$ is again a real symmetric
circulant, so it commutes with $T^{-1}$, and
\[
 D^2\log|\det T|[H,H]=-\mathop{\rm Tr}((T^{-1}H)^2)\leq0.
\]
The coordinate logarithms have strictly negative Hessian on the free
coordinates. Consequently the critical point is nondegenerate.
The critical equations are rational over $\mathbb Q(\delta)$;
clearing their nonzero denominators gives an isolated simple complex
algebraic point. Hence each $u_a$ is algebraic.

Put $R_a=\delta/u_a$ for $2\leq a<n$. The phase definitions imply
\[
 W_1^n=\delta\prod_{a=2}^{n-1}R_a,\qquad
 W_a=\frac{W_1^a}{\prod_{j=2}^{a}R_j}\quad(2\leq a<n).
\]
They make every $W_a$ algebraic; \eqref{foundation:phases} then does
the same for $F(a)$.
\end{proof}

\section{Convolution and an exact Galois symmetry}

\begin{proposition}\label{fd:frobenius}
Let $K=\mathbb Q(\delta)$ and $E=K(W_0,\ldots,W_{n-1})$.
Every phase is a unit at every prime above $2$, the extension $E/K$
is unramified there, and $E$ has a $K$-automorphism $\sigma$ satisfying
\begin{equation}\label{foundation:exact-frobenius}
 \sigma(W_a)=W_{2a}\qquad(a\in\mathbb Z_n).
\end{equation}
\end{proposition}

\begin{proof}
First, the zero Fourier mode gives a useful exact identity.
For $0<h<n$ let $r\in\{1,\ldots,n-1\}$ represent $h/2$ modulo $n$.
If $h$ is odd, $r=(n+h)/2$ and both gamma arguments in
\eqref{fd:fourier} at $k=0$ equal $B+r$.
If $h$ is even, $r=h/2$ and the arguments are $r,\Omega+r$;
their product is $(-1)^r/F(r)^2$, with this sign canceled by the
exponential factor. In both cases,
\[
 \sum_a\mathsf B_{a,h}=-B/F(r)^2.
\]
Since $F(r)^2=dW_r^2$ and
$\mathsf B_{a,h}=W_aW_{h-a}/W_h$ for $h\ne0$, this is
\begin{equation}\label{foundation:convolution}
 \sum_a W_aW_{2r-a}
 =-\delta\tau\,\frac{W_{2r}}{W_r^2}\qquad(r\in\mathbb Z_n).
\end{equation}
The case $r=0$ follows separately from
$1+(n-1)\delta=\delta(\delta-1)$ and $W_0=-1$.

Fix a valuation $v$ above $2$. Both $\delta$ and $\tau$ are units:
the reduction of $X^2-nX-1$ is $X^2+X+1$.
If $m=\min_a v(W_a)<0$, choose $r$ attaining the minimum.
Pair the summands in \eqref{foundation:convolution} under
$a\leftrightarrow2r-a$. Because $n$ is odd, the unique fixed summand
is $W_r^2$; every other pair has valuation at least $v(2)+2m>2m$.
Thus the left side has valuation $2m$, forcing
$v(W_{2r})=4m<m$, a contradiction. Every phase is integral;
the inverse norms make all phases units. Reduction now gives
\begin{equation}\label{foundation:residue-frobenius}
 \overline W_{2r}=\overline W_r^{\,4}.
\end{equation}

The completion $K_2$ is the unramified quadratic extension of
$\mathbb Q_2$. Put $\lambda=\delta\tau$ and consider the polynomial
system
\[
 \mathcal H_r(X)=X_r^2\sum_aX_aX_{2r-a}+\lambda X_{2r}=0.
\]
Its reduction is $X_r^4+X_{2r}=0$, with permutation-matrix
Jacobian. A direct form of its Hensel lift is
\[
 (\mathcal T X)_{2r}=-\lambda^{-1}X_r^2\sum_aX_aX_{2r-a}.
\]
Choose a complete unramified integer ring containing the prescribed
residue tuple; such a ring exists by the finite \'etale correspondence
for henselian rings \cite[Tag~04GK]{Stacks}.
This map preserves the tuple and improves
$X-Y\in2^s\mathcal O^n$ to
$\mathcal T X-\mathcal T Y\in2^{s+1}\mathcal O^n$ for $s\geq1$.
Indeed, writing $X=Y+e$, the convolution difference is
$2Y*e+e*e$ and the square difference has the same divisibility.
Iteration therefore gives a unique integral lift in that unramified
ring.

We also need uniqueness against ramified extensions. Normalize
$v(2)=1$, and suppose two integral tuples have the same residue,
with $s=\min_a v(X_a-Y_a)>0$. The same computation gives
\[
 \min_a v((\mathcal T X)_a-(\mathcal T Y)_a)
 \geq\min\{s+1,2s\}>s.
\]
Consequently two distinct fixed points cannot have the same residues.
The completed analytic phases are fixed points by
\eqref{foundation:convolution}; they therefore coincide with their
unramified lift. This proves unramifiedness of $E/K$ at every prime
above $2$.

At any such completion let $\varphi$ be arithmetic Frobenius over
$K_2$, with residue action $x\mapsto x^4$. Both tuples
$(\varphi(W_a))_a$ and $(W_{2a})_a$ solve the convolution system:
the first because $\varphi$ fixes $\lambda$, the second by reindexing.
Their residues agree by \eqref{foundation:residue-frobenius}.
Uniqueness proves $\varphi(W_a)=W_{2a}$. Since all these images
belong to $E$ and doubling permutes the indices, restriction gives
the automorphism \eqref{foundation:exact-frobenius}.
\end{proof}

In particular, whenever $2^m\equiv-1\pmod n$, the automorphism
$\sigma^m$ fixes $\delta$ and sends $W_a$ to $W_{-a}$ and
$\mathsf B_{a,b}$ to $\mathsf B_{-a,-b}$. Thus simultaneous negation
of indices preserves the vanishing of every polynomial residual with
coefficients in $K$.
\section{Product Fourier transforms and their contour corrections}
\label{an:products}

Write $u_p(a)=\mathsf B_{a,p}$, with subscripts in $\mathbb Z/n\mathbb Z$, and set
\[
 S_{p,q,s}(k)=\sum_{a\bmod n}u_p(a)u_q(a+s)e^{-2\pi ika/n}.
\]
For $g,h,g+h\ne0$, let
\[
 \mathcal C_{g,h}(k,l)=\sum_m\mathsf B_{m,g+h}
    \mathsf B_{g,m+k}\mathsf B_{h,m+l}
       -\tau\mathsf B_{g+l,k}\mathsf B_{h+k,l}.
\]
The order-three symmetry gives $\mathsf B_{g,a}=u_{-g}(a-g)$.
Fourier transformation in $k$, followed in the second term by
$a=k-g-l$, therefore gives the exact finite identity
\begin{equation}
 \sum_k\mathcal C_{g,h}(k,l)e^{-2\pi ijk/n}
 =e^{-2\pi ijg/n}\left[
  \widehat u_{-g}(j)S_{g+h,-h,l-h}(-j)
       -\tau e^{-2\pi ijl/n}S_{-g-l,l,g+h+l}(j)\right].
 \label{an:finite}
\end{equation}
This reindexing uses no quadratic or cubic equation.

For signed representatives $-n<q<0$, the column interpolation takes
the particularly useful form
\begin{equation}
 u_q(a)=\frac{(-1)^{q(a+1)}}{\sqrt d\,F(q)}
                      \frac{F(a)}{F(a-q)},\qquad 0\le a<n.
 \label{an:signed}
\end{equation}
There is no sample correction: at $a=0$ the inverse-norm relation gives
$-1$, and at $a=n+q$ the identity $F(n)=1/\sqrt d$ gives $-1$.
All other samples follow directly from the phase formula and the
integer sine-shift identities.

For arbitrary integer labels introduce
\[
 R_p(z)=\frac{F(z)}{F(z-p)},\qquad
 K(z)=\frac{\pi e^{-i\omega z}}{\sin\pi z}R_p(z)R_q(z+s),\qquad
 \omega\equiv\frac{2\pi k}{n}+\pi(p+q-1)\pmod{2\pi},
\]
and
\begin{align*}
 J_1(z)&=\gamma(-z)\gamma(z-p)
                   \gamma(B-z-s)\gamma(A+z+s-q),\\
 J_2(z)&=\gamma(A-z)\gamma(B+z-p)
                   \gamma(-z-s)\gamma(z+s-q).
\end{align*}
The sine-shift equations give the meromorphic identity
\begin{equation}
 K(z+n)-K(z)=-4\pi e^{-i\omega z}
 \left[(-1)^q\sin\frac{\pi p}{A}J_1(z)
       +(-1)^s\sin\frac{\pi q}{A}J_2(z)\right].
 \label{an:two-ratio}
\end{equation}
Indeed, put $M_p=R_p(z+n)/R_p(z)$ and
$M_q=R_q(z+s+n)/R_q(z+s)$ and split the multiplier into
\[
 (-1)^{p+q}M_pM_q-1
 =\bigl((-1)^pM_p-1\bigr)(-1)^qM_q+
                                      \bigl((-1)^qM_q-1\bigr).
\]
Apply the single-ratio sine identity underlying \eqref{fd:fourier}
to each bracket.  The remaining factors are
$R_q(z+s+n)=\gamma(B-z-s)\gamma(A+z+s-q)$ and
$R_p(z)=\gamma(A-z)\gamma(B+z-p)$; also
$\sin\pi(z+s)/\sin\pi z=(-1)^s$.

In the rest of this section assume
\begin{equation}
 g,h\ge1,\qquad p=g+h<n,\qquad 1\le l\le h,
 \qquad q=-h,\quad s=l-h.
 \label{an:ordered}
\end{equation}
For any labels define
\[
 c_0=\frac{(-1)^{p+q+qs}}{dF(p)F(q)},\qquad
 a_0=(-1)^q\sin\frac{\pi p}{A},\qquad
 b_0=(-1)^s\sin\frac{\pi q}{A},\qquad
 I_i=\int_{\mathcal C_i}e^{-i\omega z}J_i(z)
                                    \frac{dz}{i\sqrt\Omega},
\]
where $\mathcal C_i$ is the meromorphically continued separating
contour for its gamma factors.

\begin{lemma}\label{an:product-formula}
For the labels in \eqref{an:ordered}, at every discrete frequency,
\begin{equation}
 S_{p,q,s}(k)=-2\sqrt\Omega\,c_0(a_0I_1+b_0I_2).
 \label{an:product}
\end{equation}
The same formula holds for
\[
 (p',q',s')=(-g-l,l,g+h+l)
\]
provided its principal frequency satisfies
\begin{equation}
 |\omega|<\pi(1-g/\Omega).
 \label{an:frequency}
\end{equation}
\end{lemma}

\begin{proof}
For the first assertion let $L=\lfloor p/B\rfloor$ and choose
$\sigma=-\varepsilon$ with
$0<\varepsilon<\min(1,(L+1)B-p)$.
In $\sigma<\Re z<\sigma+n$ there are no numerator poles of $K$;
$F(z+s-q)=F(z+l)$ has no zero.  The noninteger poles are exactly
$z_j=p-jB$, $1\le j\le L$.  The shifted kernel $K(z+n)$ is
regular there.  Irrationality of $B$ excludes coincidences with
the other divisors.

To turn the vertical integrals into separating integrals one must
include the left poles $z_0=p,z_1,\ldots,z_L$ of $J_1$ and
$z_1,\ldots,z_L$ of $J_2$.  Their other left poles lie to the left
of $-l$ and their right poles are nonnegative.  The only possibly
negative opposite-parameter sums are $-p$ and $B-g-l$.
Neither is $-rA-tB$ for nonnegative integers $r,t$: since $A>n$,
such an equality would express a positive integer smaller than $n$
as a positive integer multiple of the irrational number $B$.
Thus the continued contours are nonsingular.

Write $r_{i,j}=\operatorname*{Res}_{z=z_j}e^{-i\omega z}J_i(z)$.
Equation \eqref{an:two-ratio} implies, for $j\ge1$,
\[
 \operatorname*{Res}_{z=z_j}K(z)=4\pi(a_0r_{1,j}+b_0r_{2,j}).
\]
The changes from vertical to separating integrals add respectively
\[
 \frac{2\pi}{\sqrt\Omega}\sum_{j=0}^Lr_{1,j}
 \qquad\text{and}\qquad
 \frac{2\pi}{\sqrt\Omega}\sum_{j=1}^Lr_{2,j}.
\]
The residue rectangle consequently gives
\begin{equation}
 \sum_{a=0}^{n-1}\operatorname*{Res}_{z=a}K(z)
 =-2\sqrt\Omega(a_0I_1+b_0I_2)+4\pi a_0r_{1,0}.
 \label{an:rectangle}
\end{equation}
The horizontal boundaries vanish because
$R_p(z)R_q(z+s)=O(e^{-\pi g|\Im z|/\Omega})$ uniformly on the
finite strip and $|\omega|\le\pi$.

The negative column has no patch in this sample interval; the
positive column has only its diagonal patch.  Hence
\[
 S_{p,q,s}(k)=c_0\sum_a\operatorname*{Res}_{z=a}K(z)
                    -B u_q(p+s)e^{-2\pi ikp/n}.
\]
On the other hand,
\[
 r_{1,0}=\frac{\sqrt\Omega}{2\pi}e^{-i\omega p}
                 \gamma(-p)R_q(p+s+n),\qquad
 c_0\,4\pi a_0r_{1,0}=B u_q(p+s)e^{-2\pi ikp/n}.
\]
Here one uses $2\sin(\pi p/A)\gamma(-p)=(-1)^pF(p)$ and,
at this integer sample, $R_q(p+s+n)=(-1)^qR_q(p+s)$.
This proves \eqref{an:product}.

For the primed labels put $G=g+l$, $H=g+h$, $t=H+l$ and
sample the same periodic sum on $a=-t,\ldots,n-t-1$.
The raw columns have precisely the patches $B$ at $a=-G$ and
$a=-H$, respectively.  This follows from the integer sine shifts
also when the interval contains negative multiples of $n$;
the numerator arguments avoid positive multiples of $n$.
Choose $\sigma=-t-\varepsilon$ with
\[
 0<\varepsilon<\min\{1,(\lfloor h/B\rfloor+1)B-h,
                         (\lfloor l/B\rfloor+1)B-l\}.
\]
There are no numerator poles.  The two noninteger pole families are
\[
 -G-jB\ (1\le j\le\lfloor h/B\rfloor),\qquad
 -H-jB\ (1\le j\le\lfloor l/B\rfloor).
\]
The shifted kernel is regular there.  The crossed left poles of
$J_1$ are the first family including $j=0$; those of $J_2$ are
the first family with $j\ge1$ and the second family including
$j=0$.  Its other left poles are left of $\sigma$, and all right
poles are right of $\sigma$.  All noninteger residue contributions
cancel by \eqref{an:two-ratio}.

If $h\ne l$, the two surviving integer contributions are
\[
 B u_l(h)e^{2\pi ikG/n},\qquad
 B u_{-G}(-H)e^{2\pi ikH/n},
\]
and cancel the two patches separately.  If $h=l$, some noninteger
poles of $J_2$ and $K$ are double.  The same argument uses their
full residues, since \eqref{an:two-ratio} is a meromorphic identity.
At the common integer point the first contribution uses the actual
second column, with value $-1$, whereas the second contribution
uses the raw first column, with value $1/d$.  Their sum is
\[
 B(1/d-1)e^{2\pi ikH/n}=(d^{-2}-1)e^{2\pi ikH/n},
\]
exactly the difference between the raw product $d^{-2}$ and the
actual product $1$.  Thus the coincident-diagonal correction is
included.  Finally $p'+q'=-g$, so the ratio product grows like
$O(e^{\pi g|\Im z|/\Omega})$, and the horizontal boundaries vanish
under exactly \eqref{an:frequency}.
\end{proof}

Write a principal frequency as $\omega=\pi v/n$, with
$v\equiv g+1\pmod2$, considered modulo $2n$.  Since
\[
 n(1-g/\Omega)=n-g+2g/A,\qquad 0<2g/A<2,
\]
condition \eqref{an:frequency} is equivalent to
\begin{equation}
 |v|\le n-g.
 \label{an:admitted-set}
\end{equation}
It admits $n-g+1$ frequencies and excludes $g-1$.

\section{The continuous Euler transformation}
\label{an:euler-section}

For separating-contour parameters define
\[
 I(a,b;c,e;y)=\int e^{2\pi iyz/\Omega}
      \gamma(a+z)\gamma(b+z)\gamma(c-z)\gamma(e-z)
                                               \frac{dz}{i\sqrt\Omega}.
\]
As usual, after continuation this notation includes all contour
residues.  Put $\Sigma=a+b+c+e$ and $r=(\Sigma-\Omega)/2$.
\begin{lemma}\label{an:euler}
The following meromorphic identity holds:
\begin{equation}
 I(a,b;c,e;y)=
 \frac{\gamma(a+c)\gamma(b+e)\gamma(a+e)\gamma(b+c)}
      {\gamma(\Sigma/2+y)\gamma(\Sigma/2-y)}
 I(a-r,b-r;c-r,e-r;y).
 \label{an:euler-general}
\end{equation}
\end{lemma}

\begin{proof}
Set
\begin{gather*}
 U=\Omega/4+(r+s)/2,\qquad V=\Omega/4+(r-s)/2,\\
 K_4(r,s,x,y)=\int e^{2\pi iyz/\Omega}
       \gamma(U+z)\gamma(U-z)
       \gamma(V+z+x)\gamma(V-z-x)\frac{dz}{i\sqrt\Omega},\\
 G_r(v)=\gamma(\Omega/2+r+v)\gamma(\Omega/2+r-v).
\end{gather*}
The two-gamma Fourier integral \cite[Eq.~(2.17)]{BSS}, applied
to the two pairs and combined by Fourier convolution, gives
\begin{equation}
 K_4(r,s,x,y)=G_r(s)e^{-2\pi ixy/\Omega}
                                  K_4(-r,-s,-y,x).
 \label{an:D}
\end{equation}
For precision, the transform of $\gamma(U+z)\gamma(U-z)$ is
$\gamma(2U)\gamma(\Omega/2-U+y)\gamma(\Omega/2-U-y)$.
Convolution, a translation of its variable by $y$, and reflection
of that variable yield \eqref{an:D} with the displayed phase.
Interchanging the two positive gamma parameters and translating
$z$ by $(s-x)/2$ gives the second identity
\begin{equation}
 K_4(r,s,x,y)=e^{\pi iy(s-x)/\Omega}K_4(r,x,s,y).
 \label{an:P}
\end{equation}

Call these operations $D$ and $P$.  The word $DPDPDP$ sends
\begin{align*}
 (r,s,x,y)&\longmapsto(-r,-s,-y,x)
 \longmapsto(-r,-y,-s,x)\\
 &\longmapsto(r,y,-x,-s)
 \longmapsto(r,-x,y,-s)\\
 &\longmapsto(-r,x,s,y)
 \longmapsto(-r,s,x,y).
\end{align*}
The six exponential exponents sum to zero.  The three Fourier
factors multiply to $G_r(s)G_r(x)/G_r(y)$, since
$G_{-r}(y)=1/G_r(y)$.  Consequently
\begin{equation}
 K_4(r,s,x,y)=\frac{G_r(s)G_r(x)}{G_r(y)}K_4(-r,s,x,y).
 \label{an:euler-centered}
\end{equation}

These compositions are legitimate before continuation.  For example,
in the neighborhood $|r|+|s|+|x|+|y|<\Omega/16$, all six
four-gamma kernels have positive real parts of their gamma parameters
on the imaginary axis and satisfy the strict exponential convergence
condition.  The convolution is first the ordinary absolutely
convergent Fourier convolution for imaginary $x,y$ and nearby real
$r,s$; holomorphy extends it to this neighborhood.  All contour
translations there cross no pole.  The identities then continue
meromorphically with their separating contours.  Finally centering
$I$ by $z\mapsto z+(c-a)/2$ turns \eqref{an:euler-centered} into
\eqref{an:euler-general}.
\end{proof}

\section{A four-term contour identity}
\label{an:four-term-section}

Continue to assume \eqref{an:ordered}, and put
\[
 G=g+l,\qquad R=g/2,\qquad
 X=(n+g)/2+l,\qquad Y=(n+g)/2+h.
\]
\begin{lemma}\label{an:four-term}
With continued separating contours one has
\begin{align}
0={}&\sin(\pi h/A)K_4(R,X,Y+B,Z)\notag\\
 &+(-1)^{G+1}\sin(\pi p/A)e^{-2\pi iZA/\Omega}
                                     K_4(R,X,Y-A,Z)\notag\\
 &+(-1)^p\sin(\pi G/A)e^{-\pi iZA/\Omega}
                                     K_4(R,X-A,Y,Z)\notag\\
 &-\sin(\pi l/A)e^{-\pi iZB/\Omega}K_4(R,X+B,Y,Z).
 \label{an:four-term-eq}
\end{align}
In the application $|\Re Z|<(\Omega-g)/2$.
\end{lemma}

\begin{proof}
First let $u=\Omega/4+(R+X)/2$ and
$v=\Omega/4+(R-X)/2$ be general parameters and define
\begin{gather*}
 \zeta=u-B+z,\quad \eta=u-A-z,\quad
 \theta=v-B+Y+z,\quad \nu=v-Y-z,\\
 c=2u-\Omega,\quad D=v+Y-u,\quad E=c+D,\quad
 f=u-v-A+Y.
\end{gather*}
Thus $\eta+\zeta=c$, $\eta+\theta=E$, and $\eta-\nu=f$.
In this proof abbreviate
\[
 a(t)=\sin(\pi t/A),\quad b(t)=\sin(\pi t/B),\quad
 \lambda=e^{2\pi iZB/\Omega},
\]
and set
\[
 \kappa_1=\frac{a(f-c)}{a(D)},\quad
 \kappa_2=\lambda\frac{b(E)}{b(D)},\quad
 \kappa_3=-\lambda\frac{b(c)}{b(D)},\quad
 \kappa_4=\frac{a(E-f)}{a(D)}.
\]
Let $T_i(z)$ be $e^{2\pi iZz/\Omega}$ times the product in row $i$:
\[
\begin{array}{c|l}
 i&\text{gamma product}\\\hline
 1&\gamma(\zeta)\gamma(\eta+\Omega)\gamma(\theta+B)\gamma(\nu)\\
 2&\gamma(\zeta+\Omega)\gamma(\eta)\gamma(\theta+B)\gamma(\nu)\\
 3&\gamma(\zeta+B)\gamma(\eta)\gamma(\theta+\Omega)\gamma(\nu)\\
 4&\gamma(\zeta+B)\gamma(\eta+\Omega)\gamma(\theta)\gamma(\nu).
\end{array}
\]
The primitive
\[
 P(z)=e^{2\pi iZz/\Omega}
       \gamma(\zeta)\gamma(\theta)\gamma(\eta+\Omega)\gamma(\nu+B)
\]
satisfies the exact meromorphic identity
\begin{equation}
 \sum_{i=1}^4\kappa_iT_i(z)=P(z+B)-P(z).
 \label{an:primitive}
\end{equation}
To verify it, divide by
$-8e^{2\pi iZz/\Omega}\gamma(\zeta)\gamma(\eta)
\gamma(\theta)\gamma(\nu)$.  The result is the sine identity
\begin{align*}
 &\kappa_1a(\eta)b(\eta)a(\theta)
 +\kappa_2a(\zeta)b(\zeta)a(\theta)
 +\kappa_3a(\zeta)a(\theta)b(\theta)
 +\kappa_4a(\zeta)a(\eta)b(\eta)\\
 &\hspace{15mm}=-b(\eta)
             \{\lambda a(\zeta)a(\theta)+a(\eta)a(\nu)\}.
\end{align*}
The coefficient of $\cos(\pi\eta/B)$ vanishes because
$\kappa_2b(c)+\kappa_3b(E)=0$.  The other coefficient reduces,
by sine addition, to
$\kappa_1a(\theta)+\kappa_4a(\zeta)=-a(\nu)$.

There is a nonempty initial domain where integration of
\eqref{an:primitive} has no residue term.  Take
\[
 u=\Omega/2+\epsilon,\quad v=B/2+\epsilon,\quad
 Y=A/2+\eta_0,\quad \Re z=-n/2,\qquad
 0<\epsilon<A/4,\quad0<|\eta_0|<\epsilon.
\]
On this line the real parts of $\zeta,\eta,\theta,\nu$ are
$\epsilon,\epsilon,\epsilon+\eta_0,\epsilon-\eta_0$, so
every table entry has a separating vertical contour.  The first
right poles of $P$ have distances $\Omega+\epsilon$ and
$B+\epsilon-\eta_0$ from this line, both greater than $B$.
The shift of its line to the right by $B$ therefore crosses no pole.
For $|\Re Z|<(A-4\epsilon)/2$ the horizontal boundaries decay
exponentially.  Integrating \eqref{an:primitive} gives zero.
The strict inequalities persist in an open complex neighborhood,
which supplies the domain for meromorphic continuation.

The four table integrals, with measure $dz/(i\sqrt\Omega)$, equal
\begin{gather*}
 e^{2\pi iZB/\Omega}K_4(R,X,Y+B,Z),\qquad
 e^{-2\pi iZA/\Omega}K_4(R,X,Y-A,Z),\\
 e^{-\pi iZA/\Omega}K_4(R,X-A,Y,Z),\qquad
 e^{\pi iZB/\Omega}K_4(R,X+B,Y,Z).
\end{gather*}
The respective substitutions in the centered kernels are
$z\mapsto z-B,z+A,z+A/2,z-B/2$.  Thus the integrated identity
continues meromorphically with the specified contours.

At the required parameters,
\[
 c=G-B,\quad D=h-l,\quad E=p-B,\quad f=p+l-B.
\]
If $h>l$, substitute
\[
 b(p-B)=(-1)^pa(p),\quad b(G-B)=(-1)^Ga(G),\quad
 b(h-l)=(-1)^{h-l+1}a(h-l)
\]
and multiply the identity by $e^{-2\pi iZB/\Omega}a(h-l)$.
This is \eqref{an:four-term-eq}.  If $h=l$, then $X=Y$;
\eqref{an:P} cancels the first and fourth terms and the second
and third terms separately, without division by $a(h-l)$.

Finally there is no contour pinch at this specialization.  The four
lists of opposite-parameter sums are
\[
\begin{array}{c|c}
 1&(A+G,\ B-l,\ \Omega+p,\ -h)\\
 2&(A+G,\ B-l,\ p,\ \Omega-h)\\
 3&(G,\ \Omega-l,\ A+p,\ B-h)\\
 4&(\Omega+G,\ -l,\ A+p,\ B-h).
\end{array}
\]
Their possibly negative entries cannot equal $-rA-tB$ for
$r,t\ge0$, because $A>n$ and $B$ is irrational.  These continued
values are therefore regular even when poles on the same side merge.
\end{proof}

\section{Vanishing of the admitted Fourier coefficients}
\label{an:vanishing-section}

For the negative column choose the principal
$\omega_R\equiv2\pi j/n+\pi(-g-1)\pmod{2\pi}$.
In the range $|\omega_R|<\pi(1-g/\Omega)$ one has
\begin{equation}
 \widehat u_{-g}(j)=-B e^{ig\omega_R/2}
  \gamma\left(\frac{\Omega-g}{2}+\frac{\Omega\omega_R}{2\pi}\right)
  \gamma\left(\frac{\Omega-g}{2}-\frac{\Omega\omega_R}{2\pi}\right).
 \label{an:signed-fourier}
\end{equation}
Indeed, use the single-column kernel identity with label $-g$ and
sample on $0,\ldots,n-1$.  Formula \eqref{an:signed} has no patch,
the rectangle has no noninteger pole, and
$-\epsilon+i\mathbb R$, $0<\epsilon<1$, already separates the
two gamma pole families.  The horizontal boundaries vanish in the
stated range.  The two-gamma Fourier evaluation and
$2\sin(-\pi g/A)\gamma(g)=(-1)^gF(-g)$ yield
\eqref{an:signed-fourier}.

\begin{proposition}\label{an:admitted}
Under \eqref{an:ordered}, the Fourier transform of
$\mathcal C_{g,h}(k,l)$ in $k$ vanishes at all $n-g+1$ frequencies
in \eqref{an:admitted-set}.
\end{proposition}

\begin{proof}
Put $\omega_L=-\omega_R$, $Z=\Omega\omega_L/(2\pi)$ and
\[
 D_0(Z)=\gamma((\Omega-g)/2+Z)\gamma((\Omega-g)/2-Z).
\]
The admitted frequencies exclude both endpoint ambiguities, so these
choices of principal frequencies agree with the two products in
\eqref{an:finite}.  For the left product put
\[
 r=-g/2,\qquad s_1=-\Omega/2-g/2-h,\qquad
 s_2=\Omega/2-g/2-h,\qquad y=-Z.
\]
Its two integrals are
$e^{\pi ipy/\Omega}K_4(r,s_1,X,y)$ and
$e^{\pi i(n+p)y/\Omega}K_4(r,s_2,X,y)$.
Apply \eqref{an:euler-centered}, then \eqref{an:P}, and reflect
the integration variable.  They become the first and second kernels
of \eqref{an:four-term-eq}.  The two right product integrals are
its third and fourth kernels, with centering factors
$e^{\pi ip'Z/\Omega}$ and $e^{\pi i(n+p')Z/\Omega}$.

Here are the complete normalizations.  After removing the common
factor $-2\sqrt\Omega$ from the product formulas, the four
coefficients in the bracket of \eqref{an:finite} are
\begin{equation}
 M\left(
 \sin\frac{\pi h}{A},\quad
 (-1)^{G+1}\sin\frac{\pi p}{A}e^{-2\pi iZA/\Omega},\quad
 (-1)^p\sin\frac{\pi G}{A}e^{-\pi iZA/\Omega},\quad
 -\sin\frac{\pi l}{A}e^{-\pi iZB/\Omega}\right),
 \label{an:coefficients}
\end{equation}
where
\[
 M=\frac{B(-1)^{g+h(l+1)}F(G)}{dF(l)}
                          e^{\pi iZ(A+l-g)/\Omega}.
\]
To verify the gamma part, the first Euler multiplier is
\[
 \frac{G_{-g/2}(s_1)G_{-g/2}(X)}{G_{-g/2}(-Z)}
 =\frac{F(p)F(G)}{F(h)F(l)}
       \frac{\sin(\pi h/A)}{\sin(\pi p/A)D_0(Z)}.
\]
The second multiplier divided by the first is
$(-1)^g\sin^2(\pi p/A)/\sin^2(\pi h/A)$.
These follow from reflection and
\[
 \gamma(A+a)=F(-a),\quad \gamma(B-a)=1/F(-a),\quad
 2\sin(\pi a/A)\gamma(-a)=(-1)^aF(a).
\]
The sample factors are
\[
 c_L=\frac{(-1)^{g+hl+1}F(h)}{dF(p)},\qquad
 c_R=\frac{(-1)^{lp+1}F(G)}{dF(l)}.
\]
Finally \eqref{an:signed-fourier} reads
$\widehat u_{-g}(j)=-B e^{-\pi igZ/\Omega}D_0(Z)$, and
$\omega_L\equiv-2\pi j/n+\pi(g-1)\pmod{2\pi}$
cancels the remaining integer phases.  These substitutions prove
\eqref{an:coefficients}, including all four signs and exponentials.
Lemma~\ref{an:four-term} now makes the bracket in
\eqref{an:finite} zero.
\end{proof}
\section{From admitted Fourier coefficients to all cubic equations}
\label{tail:completion}

We record the algebraic normalization needed for the completion argument.
Write $\tau=\delta-1$, where $\delta=-d^{-1}$ and
$\delta^2=n\delta+1$.  The phases constructed above satisfy
$W_0=-1$ and $W_aW_{-a}=\delta$ for $a\ne0$.  Their matrix is
\[
 \mathsf B_{a,b}=\delta^{-1}W_aW_{b-a}W_{-b}
                  +\frac{\tau^2}{\delta}\mathbf1_{a=b=0},
 \qquad \mathsf A_{a,b}=\frac{\mathsf B_{a,b}}{\tau}.
\]
Thus $\mathsf B_{0,0}=\delta-2$, and every other entry on the
two axes or the diagonal is $-1$.  The matrix $\mathsf A$ has
\begin{equation}\label{tail:linear}
 \mathsf A_{a,b}=\mathsf A_{-b,a-b},\qquad
 \mathsf A_{a,0}=\mathsf A_{0,a}=\mathbf1_{a=0}-\tau^{-1},
 \qquad \sum_a\mathsf A_{a,0}=-\delta^{-1}.
\end{equation}
All indices in this section are in $\mathbb Z/n\mathbb Z$ unless
integer inequalities are displayed.  Define
\begin{align}
 Q(g,h)&=\sum_m\mathsf A_{m+g,h}\mathsf A_{h,m}
                   -\mathbf1_{g=0}+\delta^{-1}\mathbf1_{h=0},
                  \label{tail:Q}\\
 C(g,h,k,l)&=\sum_m\mathsf A_{m,g+h}\mathsf A_{g,m+k}
                                  \mathsf A_{h,m+l}
     -\mathsf A_{g+l,k}\mathsf A_{h+k,l}
     +\delta^{-1}\mathbf1_{g=h=0},\label{tail:C}\\
 \mathcal C_{g,h}(k,l)&=\tau^3 C(g,h,k,l).
                  \label{tail:scaledC}
\end{align}
The previously established quadratic identities say $Q(g,h)=0$
for every $g,h$.  The admitted Fourier result in
Proposition~\ref{an:admitted} concerns the last normalization.

\subsection{Balanced triples and collisions}

For a cubic index quadruple put
\begin{equation}\label{tail:triples}
 \boldsymbol a=(0,k-g,l-h),\qquad
 \boldsymbol c=(g+h,-k,-l).
\end{equation}
Their six entries sum to zero.  The cross sums are
\begin{equation}\label{tail:crossmatrix}
 (a_i+c_j)_{i,j=0}^2=
 \begin{pmatrix}
 g+h&-k&-l\\
 h+k&-g&k-g-l\\
 g+l&l-h-k&-h
 \end{pmatrix}.
\end{equation}
Conversely, every balanced pair of triples has this form after
opposite translation: replace $a_i$ by $a_i-a_0$ and $c_j$ by
$c_j+a_0$, then take $k=-c_1$, $l=-c_2$, $g=-c_1-a_1$,
and $h=-c_2-a_2$.  Call a pair \emph{cross-free} if none of
its nine cross sums is zero.

\begin{lemma}\label{tail:collisions}
Every cubic with a zero in \eqref{tail:crossmatrix} follows
from the quadratic identities and the phase inverse norms.
For a cross-free pair define
\[
 S_W(\boldsymbol a,\boldsymbol c)
   =\sum_m\prod_{i=0}^2 W_{m+a_i}\prod_{j=0}^2W_{c_j-m},
 \qquad P_W(\boldsymbol a,\boldsymbol c)
   =\prod_{i,j=0}^2W_{a_i+c_j}.
\]
Then
\begin{equation}\label{tail:factor}
 \mathcal C_{g,h}(k,l)
 =\frac{W_{-g-h}W_gW_h}{\delta^5}
       \left\{\delta^2 S_W(\boldsymbol a,\boldsymbol c)
                      -\tau P_W(\boldsymbol a,\boldsymbol c)\right\}.
\end{equation}
Consequently, vanishing of a cross-free cubic is preserved
by independent permutations and opposite translation of its triples.
\end{lemma}

\begin{proof}
Choose either square root $s$ of $\delta$ and put
$b_a=W_a/s$, $\lambda=s/\tau$, and $\mu=\tau/\delta$.
These are auxiliary algebraic variables, with no complex conjugation
in their definition.  They satisfy
\[
 b_0=-s^{-1},\qquad b_ab_{-a}=1\ (a\ne0),\qquad
 \mathsf A_{u,v}=\lambda b_ub_{v-u}b_{-v}
                                  +\mu\mathbf1_{u=v=0}.
\]
Write $S_b$ and $P_b$ for the corresponding six-factor sum and
nine-factor product.  If $g,h,g+h\ne0$, direct substitution
in \eqref{tail:C} gives
\begin{equation}\label{tail:balanced}
 C(g,h,k,l)=\lambda^3b_{-g-h}b_gb_h
 \left(S_b-\lambda^{-1}P_b
       -\frac{\tau^3}{\delta^2}\,
        \mathbf1_{\{c_0,c_1,c_2\}=\{-a_0,-a_1,-a_2\}}\right).
\end{equation}
Here equality inside the indicator means equality of multisets.
Indeed, an origin correction on the right side of the cubic can
occur only at $(k,l)=(0,-g)$ or $(-h,0)$, and its normalized
coefficient is $\mu/\lambda^2=\tau^3/\delta^2$.  These are
exactly the two possible multiset matches.  Under the three
nonzero conditions on $g,h,g+h$, every such match has three
distinct row entries: a matching permutation must avoid the
three nonzero diagonal cross sums, leaving just the two
derangements, and a repeated row would also permit a forbidden match.

For a full match with distinct rows, inverse norms alone give
\[
 S_b=n-3+3\delta^{-1},\qquad
 \lambda^{-1}P_b=-\tau\delta^{-2}.
\]
Their difference is $\tau^3/\delta^2$, using
$n=\delta-\delta^{-1}$.  For a collision without a full match,
permute so that $a_0+c_0=0$.  Then
\[
 b_{m+a_0}b_{-m-a_0}
    =1+(\delta^{-1}-1)\mathbf1_{m=-a_0}.
\]
The unrestricted sum of the other four factors is, after
reindexing,
\[
 \sum_t b_tb_{H-t}b_{-t-V}b_{t+V-H},
 \qquad H=a_1+c_1,\quad V=a_2+c_1.
\]
Both $H$ and $V$ are nonzero, since otherwise the remaining
two pairs also match.  This sum is
$\lambda^{-2}Q(-V,H)$ and hence vanishes.
The surviving contribution to $S_b$ is
\[
 (\delta^{-1}-1)
       \prod_{i=1}^2b_{a_i-a_0}\prod_{j=1}^2b_{c_j+a_0}.
\]
The four interior factors of $P_b$ pair to one; thus $P_b$ is
$b_0$ times the displayed boundary product.  Since
$\lambda^{-1}b_0=\delta^{-1}-1$, this proves the collision
case of \eqref{tail:balanced}.

For completeness, the boundary cases $g=0$, $h=0$, or $g+h=0$
can be checked without dividing by the preceding prefactor.
Set $c=\tau^{-1}$.  For $g=0$, $h\ne0$, the cubic sum reduces
by $Q=0$ to
\[
 \mathsf A_{-k,h}\mathsf A_{h,l-k}-c\mathbf1_{l=0}.
\]
For $l\ne0$ the phase formula identifies this with
$\mathsf A_{l,k}\mathsf A_{h+k,l}$.  For $l=0$ use
\[
 \mathsf A_{u,h}\mathsf A_{h,u}=
 \begin{cases}c^2,&u=0\text{ or }u=h,\\
               \lambda^2,&u\ne0,h,
 \end{cases}
 \qquad (h\ne0),\qquad \lambda^2-c=c^2.
\]
The case $h=0$, $g\ne0$ follows by interchanging $(g,k)$
and $(h,l)$.  For $h=-g\ne0$ the cubic sum reduces to
\[
 \mathsf A_{g,k}\mathsf A_{-l,g}-c\mathbf1_{k=g+l}.
\]
Again the phase formula handles $k\ne g+l$, and the two
product values above handle $k=g+l$.
Finally, for $g=h=0$ the cubic sum is
\[
 \mathbf1_{k=l=0}
 -c(\mathbf1_{k=0}+\mathbf1_{l=0}+\mathbf1_{k=l})
 +3c^2-nc^3.
\]
It equals $\mathsf A_{l,k}\mathsf A_{k,l}-\delta^{-1}$:
the respective product values are $(1-c)^2$, $c^2$, and
$\lambda^2$, according as $k=l=0$, exactly one of
$k=0,l=0,k=l$ holds, or none holds.  The dimension equation
verifies all three equalities.  This proves all boundary cubics.

In the cross-free case there is no origin correction in any
factor of the cubic.  Substitution of the phase formula and
$W_aW_{-a}=\delta$ gives \eqref{tail:factor}; equivalently it
follows from \eqref{tail:balanced} using $S_W=s^6S_b$ and
$P_W=s^9P_b$.  Its prefactor is nonzero.  Both $S_W$ and $P_W$
are invariant under independent permutations, and opposite
translation is absorbed by reindexing $m$ in $S_W$.
\end{proof}

\subsection{An induction on the first index}

\begin{proposition}\label{tail:ordered}
At every odd $n\ge3$, all ordered cubics vanish:
\[
 \mathcal C_{g,h}(k,l)=0\quad
 (g,h\ge1,\ g+h<n,\ 1\le l\le h,\ k\in\mathbb Z/n\mathbb Z).
\]
\end{proposition}

\begin{proof}
Fix $n$ and induct on the integer $g$.  By
Proposition~\ref{an:admitted}, for each ordered $g,h,l$, precisely
$n-g+1$ of the $n$ Fourier coefficients of the sequence
$k\mapsto\mathcal C_{g,h}(k,l)$ are known to vanish.
The omitted frequencies form a consecutive cyclic block of
length $g-1$: the principal integers $v=n\omega_R/\pi$,
of parity $g+1$, advance by two as the Fourier index advances
by one, and the admitted condition is $|v|\le n-g$.
For $g=1$ there are no omitted frequencies, proving the base case.
For $g>1$, Fourier inversion therefore gives, with
$\zeta=\exp(2\pi i/n)$,
\begin{equation}\label{tail:polynomial}
 \mathcal C_{g,h}(k,l)=\zeta^{ak}P(\zeta^k),\qquad
                         \deg P\le g-2
\end{equation}
for an integer $a$ and a complex polynomial $P$.

We show that $k=1,\ldots,g-1$ are zeros.  A collision is
already covered by Lemma~\ref{tail:collisions}.  Otherwise
perform the following changes of cubic labels:
\begin{equation}\label{tail:descent}
 (g,h,k,l)\longmapsto
 \begin{cases}
 (k,h,g,l+g-k),&l+g-k\le h,\\
 (k,n-h-k,k+h-l,n-g-l),&l+g-k>h.
 \end{cases}
\end{equation}
These changes arise from permutations and opposite translation
of \eqref{tail:triples}.  With positions numbered $0,1,2$,
the first uses rows $(1,0,2)$ and unchanged columns; the second
uses rows $(2,0,1)$ and columns $(2,1,0)$.  In both cases
translate the first row entry to zero.  Substitution gives
exactly \eqref{tail:descent}, with all entries interpreted
modulo $n$ where necessary.

The new first index is $k<g$.  Both new quadruples are ordered.
This is immediate in the first case; in the second,
\[
 k+(n-h-k)=n-h<n,\qquad
 1\le n-g-l\le n-h-k,
\]
where the last inequality follows from $l+g-k>h$.
The induction hypothesis and the nonzero factor in
\eqref{tail:factor} therefore give each required zero of the
original cubic.  The $g-1$ points
$\zeta,\ldots,\zeta^{g-1}$ are distinct, including for
composite $n$, so \eqref{tail:polynomial} forces $P=0$.
\end{proof}

\subsection{A signed normal form}

\begin{lemma}\label{tail:normalform}
Every cross-free balanced pair of triples can be put into
the ordered range of Proposition~\ref{tail:ordered} by
independent permutations and opposite translation, possibly
after negating all six entries simultaneously.
\end{lemma}

\begin{proof}
Let $X_{ij}\in\{1,\ldots,n-1\}$ be the integer representative
of $a_i+c_j$.  Balance implies
$X_{00}+X_{11}+X_{22}\in\{n,2n\}$.  Negation replaces each
$X_{ij}$ by $n-X_{ij}$, so we may arrange that the diagonal
sum is $2n$.

There exist $i\ne j$ with $X_{ij}\ge X_{jj}$.  Otherwise
each diagonal entry would be a strict maximum in its column.
The three row labels would then be distinct.  Let $d_j>0$
be the forward cyclic gap from $a_j$ to the next row label.
Since the other entry in column $j$ is smaller and no cross
sum is zero, $X_{jj}+d_j>n$.  The gaps sum to $n$, so summing
would give $\sum_jX_{jj}>2n$, a contradiction.

Choose such $i,j$, let $r$ be the remaining position, and use
row and column order $(i,r,j)$.  Normalize the first row to
zero.  The resulting cubic indices have
\[
 g=n-X_{rr},\qquad h=n-X_{jj},\qquad l=n-X_{ij}.
\]
Thus $g,h\ge1$, $1\le l\le h$, and
$g+h=2n-X_{rr}-X_{jj}=X_{ii}<n$, as required.
\end{proof}

\section{Reconstruction and infinitude}
\label{tail:reconstruction}

Suppose $2^m\equiv-1\pmod n$ for some $m\ge1$.
The characteristic-zero Frobenius automorphism established in
Proposition~\ref{fd:frobenius} fixes $\delta$ and sends $W_a$ to $W_{2a}$.
Consequently its $m$th power sends
\[
 W_a\longmapsto W_{-a},\qquad
 \mathsf B_{a,b}\longmapsto\mathsf B_{-a,-b}.
\]
Applying it to a cubic and replacing the summation index by
its negative shows that simultaneous negation of all four
cubic indices preserves vanishing.  Lemma~\ref{tail:normalform},
Proposition~\ref{tail:ordered}, and Lemma~\ref{tail:collisions}
now prove $C(g,h,k,l)=0$ for every quadruple.

We explicitly check that these equations imply all hypotheses
of Evans--Gannon reconstruction.  In addition to
\eqref{tail:linear} and $Q=0$, their equations require the
quartic residual
\begin{align*}
 K(g,h,i,k)={}&\sum_{l,m}\mathsf A_{l,m}\mathsf A_{l+g,h}
                  \mathsf A_{h+m,l+i}\mathsf A_{i,k+m}
       -\mathsf A_{h-g,i-g}\mathbf1_{k=g}\\
 &\quad+\delta^{-1}\mathbf1_{h=0}\mathsf A_{i,k}
       +\delta^{-1}\mathsf A_{g,h}\mathbf1_{i=0}
\end{align*}
to vanish.  Order-three symmetry alone gives the polynomial identity
\begin{equation}\label{tail:quartic}
 K(g,h,i,k)=\sum_m\mathsf A_{i,k+m}
                          C(-h,h+m,g-h,i)
                  +\mathsf A_{i-h,g-h}Q(g-k,i).
\end{equation}
Indeed, replace $\mathsf A_{l+g,h}$ by
$\mathsf A_{-h,l+g-h}$ in the double sum.  The inner sum is
\[
 C(-h,h+m,g-h,i)+\mathsf A_{i-h,g-h}\mathsf A_{m+g,i}
                       -\delta^{-1}\mathbf1_{h=0}\mathbf1_{m=0}.
\]
Multiplication by $\mathsf A_{i,k+m}$ and summation over $m$
leave the quadratic sum
$Q(g-k,i)+\mathbf1_{g=k}-\delta^{-1}\mathbf1_{i=0}$.
The remaining terms cancel using
$\mathsf A_{i-h,g-h}=\mathsf A_{h-g,i-g}$ and, when $i=0$,
$\mathsf A_{-h,g-h}=\mathsf A_{g,h}$.  This proves
\eqref{tail:quartic}, hence $K=0$.

Equations \eqref{tail:linear}, \eqref{tail:Q}, and
\eqref{tail:quartic} are exactly equations (4.7)--(4.10) of
\cite{EG} with $\omega=1$ and the negative choice
$\delta=(n-\sqrt{n^2+4})/2$.  Theorem 2(a) of \cite{EG}
allows this choice.  It constructs a spherical fusion category
with simple objects $\alpha_a$ and $\alpha_a\rho$, $a\in\mathbb Z_n$,
and fusion rules
\begin{gather*}
 \alpha_a\alpha_b=\alpha_{a+b},\qquad
 \alpha_a(\alpha_b\rho)=\alpha_{a+b}\rho,
 \qquad (\alpha_b\rho)\alpha_a=\alpha_{b-a}\rho,\\
 (\alpha_a\rho)(\alpha_b\rho)
      =\alpha_{a-b}+\sum_{c\in\mathbb Z_n}\alpha_c\rho.
\end{gather*}
The categorical dimensions are $1$ and $\delta$.

The entries are algebraic by Lemma~\ref{foundation:algebraicity}.  An
embedding of $\mathbb Q(\delta)$ taking $\delta$ to
$d=(n+\sqrt{n^2+4})/2$ extends to the number field of matrix
coefficients.  Apply this embedding to the reconstruction
equations and reconstruct again.  The resulting spherical
category has categorical dimensions $1$ and $d$, which are
the Frobenius--Perron dimensions of the displayed fusion rules.
Its global dimension is therefore $n(1+d^2)$, equal to its
Frobenius--Perron dimension, so it is pseudo-unitary.
This conclusion does not assert unitarity: Theorem 2(c) of
\cite{EG} additionally requires the positively conjugated
matrix to be Hermitian, a property not established here.

Finally, $n=2^m+1$ is odd and $2^m\equiv-1\pmod n$ for every
$m\ge1$.  These orders are unbounded and the resulting
categories have distinct ranks $2n$.  Hence there are
infinitely many complex spherical cyclic Haagerup--Izumi
fusion categories, and infinitely many pseudo-unitary such
categorifications.

\begingroup
\small
\begin{thebibliography}{9}
\setlength{\itemsep}{3pt}
\bibitem{BSS}
N. M. Belousov, G. A. Sarkissian and V. P. Spiridonov,
\emph{Complex binomial theorem and pentagon identities},
arXiv:2412.07562v2 (2025), especially equations (1.4)--(1.5) and (2.17)--(2.18).
\href{https://arxiv.org/abs/2412.07562v2}{arXiv:2412.07562v2}.

\bibitem{EG}
D. E. Evans and T. Gannon,
\emph{Non-unitary fusion categories and their doubles via endomorphisms},
\emph{Advances in Mathematics} \textbf{310} (2017), 1--43,
especially Theorem 2 and equations (4.7)--(4.10).
\href{https://arxiv.org/abs/1506.03546}{arXiv:1506.03546}.

\bibitem{Stacks}
The Stacks Project Authors, \emph{The Stacks Project},
Lemma 10.153.7 (Tag 04GK), finite \'etale extensions of henselian local rings.
\href{https://stacks.math.columbia.edu/tag/04GK}{Tag 04GK}.
\end{thebibliography}
\endgroup
\end{document}
